\documentclass[11pt,showkeys]{../../shared/essay}

% Optionally override the glossary file for this paper (uncomment if needed)
% \renewcommand{\glossaryfile}{../glossaries}

\begin{document}

\title{\textit{Rooted in the Land:}\\
Redesigning GPS Through Indigenous Perspectives}

\author{Jaxen Dutta}
  \email[]{jaxen.dutta@uwaterloo.ca}
  \affiliation{David R. Cheriton School of Computer Science,\\ University of Waterloo,
  Waterloo, Ontario, Canada}

\date{February 5, 2024}

\begin{abstract}
The GPS was designed around Western cartographic conventions: standardized coordinates, abstracted overhead views, and value-neutral maps that strip landscapes of cultural meaning. Indigenous communities worldwide navigate through radically different epistemologies, grounded in relational knowledge of land, celestial bodies, ecological cues, and oral tradition. This essay examines those perspectives and argues that a genuinely inclusive redesign of GPS technology must move beyond surface-level localization to integrate four substantive changes: support for Indigenous mapping techniques and environmental navigation cues, meaningful community co-design throughout the development process, built-in cultural sensitivity protections for sacred and restricted sites, and features that promote ecological interconnectedness rather than detached bird's-eye abstraction. Two case studies, the CyberTracker project in South Africa and the Maori Maps initiative in New Zealand, demonstrate that such integration is not only possible but produces technology that is more accurate, culturally preserving, and community-empowering than any top-down alternative.
\end{abstract}

\keywords{GPS, Indigenous knowledge, participatory design, cultural sensitivity, traditional ecological knowledge, decolonial technology, navigation, cartography}

\maketitle

\section{Introduction}

The \gls{gps} has become an integral part of daily life, revolutionizing navigation and location-based services. Yet its development has predominantly been shaped by Western cartographic and engineering traditions, producing a system that reflects a particular way of knowing and representing the land: standardized coordinates, abstracted overhead views, and maps that render landscapes as neutral geometric grids. For Indigenous communities worldwide, this paradigm is neither neutral nor universal.

Indigenous navigation systems are grounded in relational, embodied, and culturally specific knowledge passed down through generations. They encompass holistic understandings of landscapes, celestial bodies, seasonal patterns, and ecological cues that form a navigation framework inseparable from cultural practice and identity. The current iteration of \gls{gps} technology not only overlooks this rich tapestry of knowledge but may actively undermine it by substituting an abstracted digital representation for the lived, relational experience of place. Understanding these perspectives is fundamental to reimagining \gls{gps} technology in a way that respects, celebrates, and integrates Indigenous wisdom.

\section{Indigenous Perspectives on Navigation and Place}

Indigenous communities share a unique connection with the land, guided by intricate traditional knowledge systems transmitted across generations. These systems encompass a holistic understanding of landscapes, celestial bodies, and environmental cues, forming a navigation framework deeply rooted in cultural practices. Where Western cartography abstracts the land into a coordinate grid, Indigenous navigation situates the traveler within a living, relational web of ecological and cultural meaning.

Furthermore, Indigenous perspectives emphasize interconnectedness, viewing humans as integral components of ecosystems rather than separate entities observing them from above. The bird's-eye view that \gls{gps} imposes, precise and efficient as it is, may inadvertently reinforce a sense of disconnection from the environment among Indigenous users. Handayani and Triyanto~\cite{handayani_integrating_2023} argue that integrating Indigenous technology into broader scientific and technological frameworks requires treating traditional knowledge not as a supplement to Western systems but as an epistemically distinct and equally valid mode of understanding the world.

This recognition carries design implications. A \gls{gps} system built for and with Indigenous communities cannot simply translate existing functionality into local languages or overlay Indigenous place names onto existing maps. It must rethink the underlying model of what navigation is for and what it means to know where you are.

\section{Four Pillars of Redesign}

\subsection{Integrating Traditional Knowledge}

The redesign of \gls{gps} technology should move beyond conventional mapping approaches and actively embrace Indigenous mapping techniques. This could involve incorporating features that allow users to navigate based on environmental cues, celestial bodies, or specific topographical features valued by Indigenous communities. By melding technological precision with traditional wisdom, the \gls{gps} system becomes a tool for cultural preservation and a living testament to Indigenous knowledge systems~\cite{handayani_integrating_2023}. Traditional ecological knowledge, including the recognition of seasonal indicators, migratory patterns, and sacred landmarks, could be encoded as navigational layers selectable by the community rather than imposed by developers.

\subsection{Community Co-Design}

Meaningful collaboration with Indigenous communities throughout the design process is not a recommendation but a foundational requirement. True co-design requires ongoing consultation, iterative user testing within communities, and a genuine commitment to incorporating feedback at every stage of development. The First Nations Technology Council~\cite{firstnations_innovation_2022} emphasizes that Indigenous-led innovation in technology is not simply a matter of representation but of sovereignty: communities must hold decision-making authority over how their knowledge is encoded, displayed, and shared. Without this structural commitment, collaboration collapses into consultation, and consultation into extraction.

\subsection{Cultural Sensitivity Protections}

The redesigned \gls{gps} technology should actively protect culturally sensitive knowledge rather than merely acknowledging its existence. The application could include features designating areas as culturally restricted, providing guidance for users to tread respectfully or avoid those areas entirely. By integrating cultural sensitivity into the technology at the architecture level, rather than as an afterthought, it becomes a tool that not only facilitates navigation but also fosters respect for Indigenous heritage and enforces community-defined boundaries on knowledge access.

\subsection{Promoting Interconnectedness}

An enhanced \gls{gps} system could transcend its role as a navigation tool by incorporating features that highlight local flora, fauna, and significant cultural sites, fostering a sense of connection between users and the environment~\cite{firstnations_innovation_2022}. This shift from the detached coordinate view to a relational, ecologically-embedded representation of place aligns the technology with Indigenous ontologies of land.

\section{Case Studies}

\subsection{CyberTracker: South Africa}

One exemplary case study is the CyberTracker project in South Africa. Indigenous trackers actively contributed to the design of a \gls{gps}-based animal tracking system, resulting in a collaboration that improved wildlife conservation efforts while ensuring the technology respected and acknowledged Indigenous tracking knowledge~\cite{cross_indigenous_2020,liebenberg_citizen_2015}. The trackers did not merely provide data: they shaped the interface, the data categories, and the representational logic of the system. This success underscores the importance of genuine co-design, recognizing Indigenous contributions as intellectual and creative authorship rather than as ethnographic material to be translated into Western technical formats.

\subsection{Maori Maps: New Zealand}

A further illuminating case study involves the Maori Maps project in New Zealand. This initiative, led by and for the Maori community, utilizes mapping technology to document and preserve Maori ancestral landscapes~\cite{oppenneer_maori_2011,scoop_maori_2021,tapsell_maorimaps_2011}. The project showcases the power of technology in cultural preservation when grounded in the principles of self-determination and community agency. The Maori Maps platform encodes whakapapa, the relational genealogical connections between people and land, directly into the navigational structure of the maps, producing a system that is simultaneously a navigation tool, a cultural archive, and an assertion of territorial sovereignty.

These case studies demonstrate that when technological innovations align with Indigenous perspectives, they become powerful tools for cultural preservation, environmental stewardship, and community empowerment.

\section{Benefits and Challenges}

The potential benefits of incorporating Indigenous perspectives into \gls{gps} technology are far-reaching. It could foster a more holistic and inclusive navigation system that acknowledges the rich knowledge systems of Indigenous communities, promotes cultural sensitivity, and deepens environmental consciousness. At the level of practical accuracy, Indigenous ecological knowledge often captures landscape dynamics that satellite imagery and administrative datasets miss entirely.

However, this transformative approach is not without its challenges. Close collaboration with Indigenous communities demands time, resources, and a commitment to addressing the power imbalances inherent in technology development. Concerns about the potential misuse or misappropriation of Indigenous knowledge necessitate robust ethical frameworks and legal protections. The co-design process must include mechanisms for communities to retain ownership of their knowledge and to withdraw consent if the technology is repurposed in ways they did not authorize.

\section{Conclusion}

The redesign of \gls{gps} technology to be more inclusive of Indigenous perspectives represents a profound shift in the paradigm of technological innovation. By acknowledging and respecting the rich knowledge systems of Indigenous communities, we pave the way for a more holistic, culturally sensitive, and interconnected technological landscape. The integration of Indigenous perspectives into \gls{gps} technology transcends mere functionality: it becomes a symbol of cultural preservation, empowerment, and environmental consciousness.

Moving forward, the design and implementation of \gls{gps} technology should prioritize ongoing collaboration with Indigenous communities, fostering relationships built on trust, respect, and shared goals. The CyberTracker and Maori Maps case studies demonstrate that this is achievable, and that the result is technology that is genuinely better: more accurate, more meaningful, and more just. As we navigate the future, our technologies should reflect the diversity of human experiences and perspectives, embracing the wisdom embedded in the landscapes and cultures of Indigenous communities around the world.

\bibliography{rooted-in-the-land}
\end{document}
